\setstretch{1.5}
\setlength{\parskip}{0.5\baselineskip}

This chapter describes the experimental methodology of each phase.
For every phase, the structure follows a consistent template: hypothesis,
method (step-by-step), variables, output files, and the research question answered.
Section~\ref{sec:design} explains the overarching design principles that govern
how the phases relate to each other.
Sections~\ref{sec:phase1} through~\ref{sec:phase6} describe each phase in full.
Table~\ref{tab:phase_summary} at the end of this chapter provides a consolidated
reference.

% ============================================================
\section{Experimental Design Principles}
\label{sec:design}
% ============================================================

\subsection{The Isolation Principle}
\label{subsec:isolation}

The defining design choice of this study is that Phases~3 through~5 are
\emph{isolated} experiments.
Each phase changes exactly one variable relative to Phase~2: the type of
knowledge injected into the system prompt.
Every other factor --- the model, the dataset, the sample set, the generation
parameters, the evaluation protocol, the analysis code --- is held fixed.
This design allows any observed CER change to be attributed unambiguously to the
knowledge injection and not to a confound.

Comparing isolated phases against Phase~2 (zero-shot) rather than against
Phase~1 (OCR baseline) is a deliberate choice.
The question of interest is not whether the LLM improves on raw OCR (Phase~2
answers that), but whether injecting a specific type of knowledge into an
already-running LLM corrector improves its performance.
Comparing to Phase~1 would confound OCR quality with LLM correction quality.

\subsection{Phase~2 as the Hub Baseline}
\label{subsec:hub}

Phase~2 (zero-shot LLM correction) serves as the comparison baseline for all
subsequent phases.
A phase that does not improve on Phase~2 is a negative result with respect to
that knowledge type: the injected knowledge does not help, and may actively
hurt, relative to the model's own language understanding.
Negative results are scientifically valuable because they rule out a hypothesis
and constrain theories of what makes LLM correction succeed.

\subsection{Reproducibility Measures}
\label{subsec:reproducibility}

All experiments use:
\begin{itemize}
  \item The same 200-sample random draw from PATS-A01 Akhbar training split.
  \item The same model: \texttt{Qwen/Qwen3-4B-Instruct-2507}.
  \item The same generation parameters: temperature 0.1, \texttt{max\_new\_tokens} 1024,
        \texttt{enable\_thinking=False}.
  \item The same evaluation protocol: diacritic-stripped CER/WER, per-sample
        then aggregated.
  \item Fixed random seeds where any sampling is involved (seed 42 throughout).
\end{itemize}
Each phase stores its inference input as a JSONL file with embedded prompts,
ensuring that the exact prompts sent to the model are permanently recorded and
can be inspected or re-run.

\subsection{Why Not Test All Combinations First}
\label{subsec:why_not_all}

With five distinct knowledge sources (confusion matrix, rules, few-shot,
self-reflective insights, RAG), the full combination space contains $2^5 = 32$
subsets.
Running all 32 on the full 200-sample dataset would require approximately
6,400 model inference calls plus the same number of analysis pipeline runs,
well beyond the available compute budget.
More importantly, running 32 uninformed combinations before understanding each
component's individual behaviour would make the results difficult to interpret:
a combination that degrades performance could not be decomposed into its
contributing factors without the isolated phase results.

The hierarchical strategy (isolated first, then best combinations in Phase~6)
extracts the maximum insight per inference call: isolated phases diagnose each
component, and Phase~6 tests whether synergies exist between the best performers.

% ============================================================
\section{Phase 1 --- Baseline and Error Taxonomy}
\label{sec:phase1}
% ============================================================

\subsection{Hypothesis}
Qaari produces systematic, categorisable errors on Arabic typewritten text.
The error profile is dominated by a small number of recurring character confusion
types that can be quantified and used to guide subsequent phases.

\subsection{Method}
\label{subsec:phase1_method}

\begin{enumerate}
  \item Load all 200 OCR predictions and ground-truth pairs for PATS-A01 Akhbar
        training split from the DataLoader.
  \item Compute per-sample CER and WER using the normalised metric
        (diacritics stripped).
        Aggregate: mean, standard deviation, median, p95.
  \item Build a character-level confusion matrix.
        For each aligned (ground-truth character, OCR character) pair in the
        minimum-edit alignment, increment the count for that pair.
        Normalise by ground-truth character frequency to obtain conditional
        confusion probabilities.
  \item Classify each substitution, insertion, and deletion into one of the
        seven taxonomic categories from Table~\ref{tab:taxonomy}
        (Chapter~3) using the heuristic rules defined there.
  \item Apply CAMeL Tools morphological analysis to each unique OCR token.
        Classify tokens as \emph{non-words} (zero valid analyses) or
        \emph{valid-but-wrong} (at least one valid analysis, but contextually
        incorrect).
  \item Write per-dataset reports including confusion matrix JSON, error
        taxonomy JSON, and a human-readable Markdown summary.
\end{enumerate}

\subsection{Output Files}
\label{subsec:phase1_outputs}
\begin{itemize}
  \item \texttt{results/phase1/\{dataset\}/metrics.json} --- CER/WER statistics
  \item \texttt{results/phase1/\{dataset\}/confusion\_matrix.json} --- character-level
        confusion counts with top-20 list
  \item \texttt{results/phase1/\{dataset\}/error\_taxonomy.json} --- error counts
        and percentages by category with examples
  \item \texttt{results/phase1/\{dataset\}/morphological\_analysis.json} --- word
        validity statistics
  \item \texttt{results/phase1/\{dataset\}/error\_examples.json} --- up to 5 sample
        OCR/GT pairs per error category
  \item \texttt{results/phase1/\{dataset\}/report.md} --- human-readable summary
\end{itemize}

\noindent
\textbf{Research question answered:} RQ1.

% ============================================================
\section{Phase 2 --- Zero-Shot LLM Correction}
\label{sec:phase2}
% ============================================================

\subsection{Hypothesis}
A zero-shot instruction-tuned LLM can improve Arabic OCR output without
task-specific guidance, by exploiting its pre-trained Arabic language knowledge
to recognise and correct implausible character sequences.

\subsection{Method}
\label{subsec:phase2_method}

\begin{enumerate}
  \item For each sample in the 200-sample evaluation set, construct the
        conservative zero-shot prompt using
        \texttt{PromptBuilder.build\_zero\_shot(ocr\_text, version="v2")}.
        The system message tells the model that the text may already be correct,
        that it should return it unchanged if no clear error is found, and that
        it must not change what is already correct.

  \item Write all prompts to
        \texttt{results/phase2/inference\_input.jsonl}, one JSON object per
        sample, containing fields: \texttt{sample\_id}, \texttt{dataset},
        \texttt{ocr\_text}, \texttt{gt\_text}, \texttt{prompt},
        \texttt{prompt\_version} (= \texttt{p2v2}).

  \item Run inference on Kaggle (GPU) using \texttt{scripts/infer.py}.
        For each sample, the model generates the corrected text.
        Results are written to \texttt{results/phase2/corrections.jsonl}.

  \item Compute post-correction CER and WER for each sample.
        Compare to Phase~1 OCR baseline.
        Compute error changes: for each sample, count errors fixed (present
        in OCR but absent from correction) and errors introduced (absent from
        OCR but present in correction).

  \item Write comparison report to
        \texttt{results/phase2/\{dataset\}/comparison\_vs\_phase1.json}.
\end{enumerate}

\subsection{Prompt Design Rationale}
\label{subsec:phase2_rationale}

Two prompt versions were developed:
\begin{itemize}
  \item \textbf{v1 (aggressive)}: ``You are an Arabic text corrector. Your task
        is to correct OCR errors in Arabic text. Return only the corrected text
        without explanation.'' (translated from Arabic)
  \item \textbf{v2 (conservative)}: ``You are an Arabic text proofreader.
        The following text is OCR output and may be fully correct or contain few
        errors. If the text is correct, return it exactly as-is. If you find only
        a clear error, correct it. Do not change what is correct. Return only the
        text without explanation.'' (translated from Arabic)
\end{itemize}
Version v1 was piloted but discarded after it produced worse aggregate CER than
the OCR baseline on clean samples: the aggressive framing caused the model to
rewrite text it should have preserved.
Version v2 is used in all experiments reported in this thesis.

\subsection{Output Files}
\begin{itemize}
  \item \texttt{results/phase2/\{dataset\}/metrics.json}
  \item \texttt{results/phase2/\{dataset\}/comparison\_vs\_phase1.json}
  \item \texttt{results/phase2/\{dataset\}/error\_changes.json}
\end{itemize}

\noindent
\textbf{Research question answered:} RQ2.
\textbf{Role in the study:} Phase~2 output serves as the comparison baseline
for all subsequent phases.

% ============================================================
\section{Phase 3 --- OCR-Aware Prompting}
\label{sec:phase3}
% ============================================================

\subsection{Hypothesis}
Telling the LLM which specific character confusions Qaari makes will direct its
attention to the positions most likely to contain errors and improve correction
quality for those error types.

\subsection{Method}
\label{subsec:phase3_method}

\begin{enumerate}
  \item Load the Phase~1 confusion matrix from
        \texttt{results/phase1/PATS-A01-Akhbar-train/confusion\_matrix.json}
        using \texttt{ConfusionMatrixLoader}.
  \item Extract the top-10 confusion pairs by substitution count.
        For the PATS-A01 Akhbar training split, the dominant confusions are
        detailed in Chapter~5.
  \item Format the confusion pairs as an Arabic-language bullet list using
        \texttt{ConfusionMatrixLoader.format\_for\_prompt()}.
        Each bullet item reads (translated): ``This system confuses [correct
        character] with [OCR character].''
  \item Build the OCR-aware prompt using
        \texttt{PromptBuilder.build\_ocr\_aware(ocr\_text, confusion\_context)}.
        The system message extends the zero-shot message with the confusion list
        and instructs the model to pay particular attention to these pairs.
  \item Export, infer, and analyze following the same three-stage pipeline as
        Phase~2.
\end{enumerate}

\subsection{Variables}
\begin{itemize}
  \item \textbf{Changed from Phase~2:} Confusion list injected into system prompt.
  \item \textbf{Fixed:} Model, dataset, generation parameters, evaluation protocol.
  \item \textbf{Number of confusion pairs:} 10 (configured as \texttt{phase3.top\_n}).
\end{itemize}

\subsection{Output Files}
\begin{itemize}
  \item \texttt{results/phase3/\{dataset\}/metrics.json}
  \item \texttt{results/phase3/\{dataset\}/comparison\_vs\_phase2.json}
  \item \texttt{results/phase3/\{dataset\}/confusion\_impact.json} ---
        per-confusion-pair analysis of whether the injection helped or hurt
        for samples containing that confusion
\end{itemize}

\noindent
\textbf{Research question answered:} RQ3.

% ============================================================
\section{Phase 4A --- Rule-Augmented Prompting}
\label{sec:phase4a}
% ============================================================

\subsection{Hypothesis}
Explicit Arabic orthographic rules provide grounded, interpretable guidance that
improves correction for rule-governed phenomena such as hamza placement, taa
marbuta, and alef maqsura.

\subsection{Method}
\label{subsec:phase4a_method}

\begin{enumerate}
  \item Load rules from \texttt{data/rules/} using \texttt{RulesLoader}.
        Rules are selected from the orthography, morphology, and syntax JSONL
        files without a category restriction (\texttt{categories: null} in config).
  \item Format rules as a numbered Arabic instruction list using
        \texttt{RulesLoader.format\_for\_prompt()}.
        The format is: ``1. [rule text in Arabic]'' repeated for each rule.
  \item Build the rule-augmented prompt using
        \texttt{PromptBuilder.build\_rule\_augmented(ocr\_text, rules\_context)}.
        The system message instructs the model to observe these rules during
        correction.
  \item Export, infer, and analyze as in Phase~2.
\end{enumerate}

\subsection{Variables}
\begin{itemize}
  \item \textbf{Changed from Phase~2:} Arabic orthographic rules injected into
        system prompt.
  \item \textbf{Fixed:} Model, dataset, generation parameters, evaluation protocol.
  \item \textbf{Rule categories:} All available (orthography, morphology, syntax).
\end{itemize}

\subsection{Output Files}
\begin{itemize}
  \item \texttt{results/phase4a/\{dataset\}/metrics.json}
  \item \texttt{results/phase4a/\{dataset\}/comparison\_vs\_phase2.json}
  \item \texttt{results/phase4a/\{dataset\}/error\_changes.json}
\end{itemize}

\noindent
\textbf{Research question answered:} RQ4.

% ============================================================
\section{Phase 4B --- Few-Shot Examples}
\label{sec:phase4b}
% ============================================================

\subsection{Hypothesis}
Providing the model with demonstration pairs of erroneous input and correct output
from the QALB corpus will give it a concrete template for what Arabic error
correction looks like, reducing aggressive rewrites and improving targeted fixes.

\subsection{Method}
\label{subsec:phase4b_method}

\begin{enumerate}
  \item Load QALB-2014 training split using \texttt{QALBLoader} from
        \texttt{data/QALB-0.9.1-Dec03-2021-SharedTasks/}.
  \item Apply filters: length 10--300 characters, maximum 15 words changed,
        OCR-relevance (dot substitutions, alef/hamza variants, similar-shape
        confusions).
        After filtering: 528 OCR-relevant pairs available.
  \item Select 5 pairs using the diversity strategy (\texttt{selection: "diverse"},
        seed 42), ensuring multiple error categories are represented among the
        5 examples.
  \item Format examples as an Arabic inline list:
        ``Input: [erroneous text] / Correction: [corrected text]'' for each pair.
  \item Build the few-shot prompt using
        \texttt{PromptBuilder.build\_few\_shot(ocr\_text, examples\_context)}.
  \item Export, infer, and analyze as in Phase~2.
\end{enumerate}

\subsection{Transfer Gap Discussion}
\label{subsec:phase4b_transfer}

The QALB corpus contains errors introduced by human writers: phonological
substitutions (writing a word as it sounds in a particular dialect), agreement
errors (wrong gender or case marking), and punctuation errors.
OCR errors are fundamentally different in kind: they are visual confusion errors
caused by similar glyph shapes.
The few-shot examples may teach the model to be attentive to Arabic spelling in
general, but the specific error patterns they demonstrate do not closely match
the error patterns produced by Qaari.
This \emph{transfer gap} is the primary structural limitation of Phase~4B and
is examined quantitatively in Chapter~5.

\subsection{Output Files}
\begin{itemize}
  \item \texttt{results/phase4b/\{dataset\}/metrics.json}
  \item \texttt{results/phase4b/\{dataset\}/comparison\_vs\_phase2.json}
  \item \texttt{results/phase4b/\{dataset\}/error\_changes.json}
\end{itemize}

\noindent
\textbf{Research question answered:} RQ5.

% ============================================================
\section{Phase 4C --- CAMeL Morphological Post-Processing}
\label{sec:phase4c}
% ============================================================

\subsection{Hypothesis}
Morphological post-processing can detect LLM hallucinations --- cases where the
LLM introduces morphologically invalid Arabic words --- and revert those words
to the (valid) OCR original, reducing the error introduction rate without
affecting correctly corrected words.

\subsection{Method}
\label{subsec:phase4c_method}

Phase~4C does not generate new prompts or run a new inference call.
It applies the revert strategy as a post-processing step to the Phase~2
corrections.

\begin{enumerate}
  \item Load Phase~2 corrections from
        \texttt{results/phase2/\{dataset\}/corrections.jsonl}.
  \item Initialise \texttt{MorphAnalyzer} with the MSA database
        \texttt{calima-msa-r13}.
        \textbf{Note:} CAMeL Tools requires a separate installation step:
        \texttt{pip install camel-tools} followed by
        \texttt{camel\_data -i morphology-db-msa-r13}.
  \item For each sample, apply \texttt{WordValidator.validate\_correction(ocr\_text, llm\_text)}:
        \begin{enumerate}
          \item Tokenise both the OCR text and the LLM correction into Arabic
                word tokens.
          \item For each word pair (ocr\_word, llm\_word) that differs:
                \begin{itemize}
                  \item Run \texttt{MorphAnalyzer.analyse(llm\_word)};
                        if result has zero analyses: LLM word is invalid.
                  \item Run \texttt{MorphAnalyzer.analyse(ocr\_word)};
                        if result has at least one analysis: OCR word is valid.
                  \item If both conditions hold: revert to \texttt{ocr\_word}.
                  \item Otherwise: keep the LLM correction.
                \end{itemize}
          \item Record: words reverted, words kept, revert rate.
        \end{enumerate}
  \item Compute post-validation CER/WER on the final text.
        Compare to both Phase~2 (LLM only) and Phase~1 (OCR baseline).
  \item Write validation statistics to
        \texttt{results/phase4c/\{dataset\}/validation\_stats.json}.
\end{enumerate}

\subsection{Revert Strategy Rationale}
\label{subsec:phase4c_rationale}

The revert strategy is deliberately conservative.
It can only revert a word to the OCR original; it cannot introduce new words.
The maximum harm it can do is revert a correctly-improved LLM word back to a
wrong OCR word --- but only if the LLM word happened to be morphologically
invalid despite being the correct intended word (which is rare for genuine
corrections).
The maximum benefit is preventing a hallucinated, morphologically invalid LLM
word from degrading the text when the OCR had a valid form at that position.

This one-sided risk profile means Phase~4C can be deployed as a defensive layer
on top of any LLM corrector with no risk of making performance worse than the
raw OCR baseline.

\subsection{Output Files}
\begin{itemize}
  \item \texttt{results/phase4c/\{dataset\}/metrics.json}
  \item \texttt{results/phase4c/\{dataset\}/comparison\_vs\_phase2.json}
  \item \texttt{results/phase4c/\{dataset\}/validation\_stats.json} ---
        per-dataset revert counts, rates, and list of reverted words
\end{itemize}

\noindent
\textbf{Research question answered:} RQ6.

% ============================================================
\section{Phase 4D --- Self-Reflective Prompting}
\label{sec:phase4d}
% ============================================================

\subsection{Hypothesis}
Feeding the model statistical summaries of its own systematic errors --- derived
from its corrections on the training split where ground truth is available ---
will cause it to apply targeted corrections more carefully and reduce its
error introduction rate on the validation split.

\subsection{Three-Stage Method}
\label{subsec:phase4d_method}

Phase~4D is itself a three-stage process that must be executed in order.

\subsubsection{Stage~1: Training-Split Analysis (analyze-train mode)}

\begin{enumerate}
  \item Load Phase~2 corrections on the training split
        (\texttt{PATS-A01-Akhbar-train}) from
        \texttt{results/phase2/PATS-A01-Akhbar-train/corrections.jsonl}.
  \item For each sample, compare the Phase~2 LLM corrections to ground truth
        using \texttt{LLMErrorAnalyzer.analyse\_sample()}.
        For every character alignment between LLM output and ground truth:
        \begin{itemize}
          \item If the OCR had an error here and the LLM fixed it: increment
                \texttt{fixed} for that error type.
          \item If the OCR was correct here and the LLM changed it incorrectly:
                increment \texttt{introduced} for that error type.
        \end{itemize}
  \item Aggregate across all training samples to obtain per-error-type statistics:
        \begin{itemize}
          \item \texttt{fix\_rate} = \texttt{fixed} / \texttt{baseline} errors of
                that type (fraction of existing errors the LLM corrected).
          \item \texttt{introduction\_rate} = \texttt{introduced} / ground-truth
                tokens of that type (fraction of correct tokens incorrectly changed).
        \end{itemize}
  \item Write aggregated statistics to
        \texttt{results/phase4d/insights/PATS-A01\_insights.json}.
\end{enumerate}

\subsubsection{Insights Derived from Training Data}

Table~\ref{tab:phase4d_insights} reports the statistics extracted from the
200-sample PATS-A01 Akhbar training split.

\begin{table}[ht]
\caption{Phase~4D training-split analysis results for PATS-A01 Akhbar.
         The LLM is Qwen3-4B-Instruct-2507 under the zero-shot v2 prompt
         (Phase~2 setting). Fix rate = fraction of baseline errors of that type
         that the LLM corrected. Introduction rate = fraction of correct ground-truth
         tokens of that type that the LLM incorrectly changed.
         $n < 10$ entries are marked with $\dagger$ (too few samples for reliable
         estimates).}
\label{tab:phase4d_insights}
\centering
\renewcommand{\arraystretch}{1.3}
\begin{tabular}{lrrrrr}
\toprule
\textbf{Error type} & \textbf{Baseline} & \textbf{Fixed} & \textbf{Introduced}
  & \textbf{Fix rate} & \textbf{Intro.\ rate} \\
\midrule
Insertion       & 1,964 & 1,912 & 356  & 97.35\% & 18.13\% \\
Other subst.    & 49    & 8     & 80   & 16.33\% & 163.27\% \\
Deletion        & 32    & 1     & 68   & 3.12\%  & 212.50\% \\
Dot confusion   & 14    & 1     & 2    & 7.14\%  & 14.29\% \\
Alef maqsura    & 3$^\dagger$ & 1 & 3 & 33.33\% & 100.00\% \\
Similar shape   & 1$^\dagger$ & 1 & 8 & 100.00\% & 800.00\% \\
Hamza           & 0     & ---   & 0    & ---     & ---     \\
Taa marbuta     & 0     & ---   & 0    & ---     & ---     \\
\midrule
\textbf{Overall} & \textbf{2,063} & \textbf{1,924} & \textbf{517}
  & \textbf{93.26\%} & \textbf{25.06\%} \\
\bottomrule
\multicolumn{6}{l}{\footnotesize $\dagger$ Fewer than 10 baseline samples; estimate unreliable.}
\end{tabular}
\end{table}

\noindent
The table reveals two important patterns.
First, the insertion fix rate of 97.35\% shows that the LLM is highly effective
at collapsing the Qaari repetition artifact: almost all spurious character
insertions (which constitute 95.2\% of all baseline errors) are removed.
Second, the overall introduction rate of 25.06\% --- one in four previously
correct tokens is incorrectly changed by the LLM --- represents a systematic
over-correction bias that affects all error categories.
The deletion and other-substitution introduction rates exceed 100\% because the
LLM introduces more deletions and substitutions of those types than existed in
the original OCR output.

\subsubsection{Stage~2: Export (export mode)}

\begin{enumerate}
  \item Load the aggregated insights from
        \texttt{results/phase4d/insights/PATS-A01\_insights.json}
        using \texttt{LLMInsightsLoader}.
  \item Apply the insight filtering logic:
        \begin{itemize}
          \item Error types with \texttt{fix\_rate} $> 0.6$: reported as
                strengths in the diagnostic.
          \item Error types with \texttt{fix\_rate} $< 0.4$ and
                \texttt{baseline} $\geq 10$: reported as weaknesses.
          \item Error types with \texttt{introduction\_rate} $> 0.05$ and
                \texttt{baseline} $\geq 10$: reported as over-correction risks.
        \end{itemize}
  \item Convert the filtered statistics to Arabic-language diagnostic sentences.
        The resulting insights text injected into the Phase~4D validation prompt
        is provided in full in Appendix~B.
        Its content identifies insertions as a strength, deletions and other
        substitutions as weaknesses, and the high overall introduction rate as
        an over-correction warning.
  \item Build the self-reflective prompt for each validation-split sample using
        \texttt{PromptBuilder.build\_self\_reflective(ocr\_text, insights\_context)}.
  \item Write to \texttt{results/phase4d/inference\_input.jsonl}.
\end{enumerate}

\subsubsection{Stage~3: Analyze (analyze mode)}

\begin{enumerate}
  \item Run inference on Kaggle GPU using \texttt{scripts/infer.py}.
  \item Load \texttt{results/phase4d/corrections.jsonl}.
  \item Compute CER/WER; compare to Phase~2.
  \item Write reports to \texttt{results/phase4d/\{dataset\}/}.
\end{enumerate}

\subsection{Output Files}
\begin{itemize}
  \item \texttt{results/phase4d/insights/PATS-A01\_insights.json} ---
        training-split statistics
  \item \texttt{results/phase4d/\{dataset\}/metrics.json}
  \item \texttt{results/phase4d/\{dataset\}/comparison\_vs\_phase2.json}
  \item \texttt{results/phase4d/\{dataset\}/error\_changes.json}
\end{itemize}

\noindent
\textbf{Research question answered:} RQ7.

% ============================================================
\section{Phase 5 --- Retrieval-Augmented Generation}
\label{sec:phase5}
% ============================================================

\subsection{Hypothesis}
Providing the model with semantically similar passages from a large Arabic corpus
(OpenITI) supplies lexical and syntactic grounding that helps it determine the
correct form of ambiguous OCR sequences.

\subsection{Method}
\label{subsec:phase5_method}

Phase~5 is itself a three-mode pipeline (build, export, analyze).

\subsubsection{Build Mode}

\begin{enumerate}
  \item Load OpenITI corpus inventory from \texttt{data/OpenITI/} using
        \texttt{OpenITILoader}.
  \item Extract up to 200,000 sentences from primary-status texts, applying
        filters: minimum 30 characters, maximum 300 characters.
        Sentence splitting is performed at punctuation boundaries (full stop,
        question mark, exclamation, Arabic semicolon U+061B).
  \item Embed all sentences using
        \texttt{paraphrase-multilingual-MiniLM-L12-v2}
        in batches of 256 on CPU.
  \item Build a FAISS \texttt{FlatIP} index from the embeddings.
        Save index and sentence list to disk.
\end{enumerate}

\subsubsection{Export Mode}

\begin{enumerate}
  \item Load the FAISS index using \texttt{RAGRetriever}.
  \item For each of the 200 evaluation samples, embed the OCR text as a query.
  \item Retrieve the top-3 most similar sentences by inner-product similarity.
  \item Format the retrieved sentences as a numbered Arabic list.
  \item Build the RAG prompt using
        \texttt{PromptBuilder.build\_rag(ocr\_text, retrieval\_context)}.
  \item Write to \texttt{results/phase5/inference\_input.jsonl}, embedding the
        retrieved text in each JSONL record so no index access is needed at
        inference time.
\end{enumerate}

\subsubsection{Analyze Mode}

\begin{enumerate}
  \item Run inference on Kaggle GPU using \texttt{scripts/infer.py}.
  \item Load \texttt{results/phase5/corrections.jsonl}.
  \item Compute CER/WER; compare to Phase~2.
  \item Compute retrieval quality metrics from the embedded retrieval scores:
        average top-1 inner-product score, average top-$k$ score,
        distribution of scores across quality bands.
        Write to \texttt{results/phase5/\{dataset\}/retrieval\_analysis.json}.
\end{enumerate}

\subsection{Retrieval Quality}
\label{subsec:phase5_retrieval_quality}

The inner-product similarity scores between the OCR queries and the retrieved
OpenITI sentences are high in absolute terms: the average top-1 score across
200 samples is 0.805 and the average top-3 score is 0.793.
However, these scores reflect embedding-space proximity as computed by
\texttt{MiniLM-L12-v2}, not semantic relevance to the correction task.
High scores indicate that the retrieved sentences are ``similar'' in embedding
space to the query, but this similarity may be surface-level (both texts are
Arabic and share common function words) rather than meaningful semantic overlap
between the retrieved classical text and the modern newspaper query.
Chapter~5 connects the retrieval score distribution to the observed CER outcome.

\subsection{Output Files}
\begin{itemize}
  \item \texttt{results/phase5/\{dataset\}/metrics.json}
  \item \texttt{results/phase5/\{dataset\}/comparison\_vs\_phase2.json}
  \item \texttt{results/phase5/\{dataset\}/retrieval\_analysis.json} ---
        per-sample retrieval scores and aggregate statistics
\end{itemize}

\noindent
\textbf{Research question answered:} RQ8.

% ============================================================
\section{Phase 6 --- Combinations and Ablation}
\label{sec:phase6}
% ============================================================

\subsection{Purpose}
\label{subsec:phase6_purpose}

Phases 3--5 establish which knowledge types help and which hurt in isolation.
Phase~6 asks two follow-up questions:
(a) can two knowledge types combine synergistically to produce a result better
than either individually?
(b) in the best combined system, what is each component's quantitative
contribution?

\subsection{Combinations Tested}
\label{subsec:phase6_combos}

The 12 combinations are selected based on the isolated phase results.
The guiding principle is that combinations involving components that performed
similarly to Phase~2 in isolation are more likely to combine beneficially than
combinations involving highly disruptive components.
Table~\ref{tab:phase6_combos} lists all 12 combinations.

\begin{table}[ht]
\caption{Phase~6 combination inventory. Each combination is identified by a
         unique ID. Components: C = confusion matrix, R = rules, E = examples
         (few-shot), S = self-reflective insights, G = RAG, M = CAMeL
         post-processing.}
\label{tab:phase6_combos}
\centering
\renewcommand{\arraystretch}{1.3}
\begin{tabular}{lll}
\toprule
\textbf{Combo ID} & \textbf{Components} & \textbf{Rationale} \\
\midrule
\texttt{conf\_rules}     & C + R   & Confusion + rules: both target character-level errors \\
\texttt{conf\_few}       & C + E   & Confusion + examples: data-driven + example-driven \\
\texttt{rules\_few}      & R + E   & Rules + examples: best two augmented phases \\
\texttt{few\_self}       & E + S   & Examples + self-reflective: both are evidence-based \\
\texttt{conf\_rules\_few}& C + R + E & Top-3 triple \\
\texttt{rules\_few\_self}& R + E + S & Alternative triple \\
\texttt{full\_no\_rag}   & C + R + E + S & Full system minus RAG (best guess given phase results) \\
\texttt{full\_system}    & C + R + E + S + G & All five knowledge types \\
\texttt{conf\_few\_camel}& C + E + M & Pair + CAMeL post-processing \\
\texttt{rules\_few\_camel}& R + E + M & Best pair + CAMeL \\
\texttt{few\_self\_d4d}  & E + S (Phase~4D) & Self-reflective as add-on \\
\texttt{full\_camel}     & C + R + E + S + M & Full system with CAMeL, no RAG \\
\bottomrule
\end{tabular}
\end{table}

\subsection{Ablation Design}
\label{subsec:phase6_ablation}

The ablation study starts from the best-performing combination identified in
Phase~6 and removes one component at a time.
The CER delta when a component is removed measures that component's individual
contribution to the combined system:

\begin{equation}\label{eq:ablation}
  \Delta_{\text{CER}}^{(k)} = \text{CER}_{\text{full} \setminus k} - \text{CER}_{\text{full}}
\end{equation}

\noindent
A positive $\Delta_{\text{CER}}^{(k)}$ means removing component $k$ increases
CER --- the component was helping.
A negative $\Delta_{\text{CER}}^{(k)}$ means removing it decreases CER ---
the component was hurting.
This design answers RQ9: which combination is optimal and what does each element
contribute?

\subsection{Statistical Testing}
\label{subsec:phase6_stats}

For Phase~6 results, paired t-tests are used to assess whether observed CER
differences between the best combination and Phase~2 are statistically
significant.
The Bonferroni correction is applied to control the family-wise error rate across
12 combination comparisons at the $\alpha = 0.05$ level.
Bootstrap confidence intervals (1,000 iterations) are also computed for the
best-performing combination.
These tests are implemented in
\texttt{src/analysis/stats\_tester.py::StatsTester}.

\subsection{Output Files}
\begin{itemize}
  \item \texttt{results/phase6/\{combo\_id\}/\{dataset\}/metrics.json}
  \item \texttt{results/phase6/\{combo\_id\}/\{dataset\}/comparison\_vs\_phase2.json}
  \item \texttt{results/phase6/summary.json} --- all combinations ranked by CER
  \item \texttt{results/phase6/ablation.json} --- ablation delta table
  \item \texttt{results/phase6/paper\_tables.md} --- formatted tables for publication
\end{itemize}

\noindent
\textbf{Research question answered:} RQ9.

% ============================================================
\section{Phase Summary}
\label{sec:phase_summary}
% ============================================================

Table~\ref{tab:phase_summary} provides a consolidated reference for all phases.

\begin{table}[ht]
\caption{Summary of all experimental phases. ``Comparison'' is the phase each
         experiment is evaluated against.}
\label{tab:phase_summary}
\centering
\renewcommand{\arraystretch}{1.3}
\begin{tabularx}{\linewidth}{lXll}
\toprule
\textbf{Phase} & \textbf{Knowledge injected} & \textbf{Comparison} & \textbf{RQ} \\
\midrule
Phase~1 & None (OCR analysis) & N/A          & RQ1 \\
Phase~2 & None (zero-shot)    & Phase~1      & RQ2 \\
Phase~3 & Confusion matrix (top-10 pairs) & Phase~2  & RQ3 \\
Phase~4A & Arabic orthographic rules & Phase~2  & RQ4 \\
Phase~4B & 5 QALB few-shot examples  & Phase~2  & RQ5 \\
Phase~4C & CAMeL morphological post-processing & Phase~2 & RQ6 \\
Phase~4D & LLM training-split error statistics & Phase~2 & RQ7 \\
Phase~5  & Top-3 OpenITI RAG passages & Phase~2  & RQ8 \\
Phase~6  & Combinations + ablation    & Phase~2 + best combo & RQ9 \\
\bottomrule
\end{tabularx}
\end{table}
